%! TEX root = master.tex
\section*{Preface}

Several mathematicians led me to this subject. The first is Vladimir Arnold,
whose work I met while studying abstract algebra in my junior year of
college. The second is Eric Weinstein, and the third is Roger Penrose.

My interest began when I became aware of the stagnation in fundamental
physics. Vast sums have gone into the search for a quantum theory of
gravity, and, as I see it, we are hardly closer to one than we were sixty
years ago. To find out where things went south, I first had to understand
which way is north. That starts with a bridge between mathematics and
classical mechanics, and these notes are my attempt to cross it. I hope they
help you as they have helped me.

The notes follow Arnold's \emph{Mathematical Methods of Classical
Mechanics} \cite{arnold}. The framework and the principles are his; the
expansions and worked details are mine.

\chapter{1}{First Principles}

Mechanics rests on two facts taken from experiment: Galileo's principle of
relativity and Newton's principle of determinacy \cite[ch.~1]{arnold}.
Before stating them we need the stage on which they act.

\section{Space and time}\label{sec:space-time}

Space is three-dimensional and Euclidean. Time is one-dimensional. The
second word needs a definition.

\begin{definition}[Euclidean space]\label{def:euclidean}
  The scalar product on $\R^n$ is
  $(\mathbf x, \mathbf y) = x_1 y_1 + \dots + x_n y_n$. It gives the distance
  \[
    \rho(\mathbf x, \mathbf y) = \sqrt{(\mathbf x - \mathbf y, \mathbf x - \mathbf y)}.
  \]
  $\R^n$ with this distance is \emph{Euclidean space} $E^n$. A space is
  \emph{Euclidean} if its distances behave like these, so that lengths and
  angles are measured as in school geometry.
\end{definition}

\section{The principle of relativity}\label{sec:relativity}

\begin{definition}[Galileo's principle of relativity]\label{def:relativity}
  There exist coordinate systems, called \emph{inertial}, with two
  properties.
  \begin{enumerate}
    \item All the laws of nature are the same, at every moment of time, in
      every inertial coordinate system.
    \item Every coordinate system in uniform rectilinear motion relative to
      an inertial one is itself inertial.
  \end{enumerate}
\end{definition}

\begin{example}\label{ex:train}
  Suppose a coordinate system attached to the earth is inertial. Then an
  experimenter in a train moving uniformly in a straight line relative to
  the earth cannot detect the train's motion by any experiment carried out
  inside the carriage.
\end{example}

\begin{remark}
  In fact a coordinate system attached to the earth is only approximately
  inertial: the earth rotates and moves on a curved orbit. The principle is
  an idealisation. Our equations describe real motions only approximately,
  so bodies moving on the earth cannot strictly be called inertial.
\end{remark}

\section{The principle of determinacy}\label{sec:determinacy}

\begin{definition}[Newton's principle of determinacy]\label{def:determinacy}
  The initial state of a mechanical system, meaning all positions and all
  velocities of its points at one moment, uniquely determines all of its
  motion.
\end{definition}

The principle says that accelerations are functions of positions,
velocities and time. Writing this down gives the equation of motion.

%! TEX root = ../master.tex
\begin{theorem}[Newton's equation]\label{thm:newton-equation}
  Let a mechanical system of $n$ points have configuration
  $\mathbf x(t) \in \R^N$, where $N = 3n$. By the determinacy principle
  (\cref{def:determinacy}) there is a function
  $\mathbf F : \R^N \times \R^N \times \R \to \R^N$ such that every motion of
  the system satisfies
  \[
    \ddot{\mathbf x} = \mathbf F(\mathbf x, \dot{\mathbf x}, t).
  \]
  Conversely, if $\mathbf F$ is smooth, then a solution is uniquely determined by
  $\mathbf x(t_0)$ and $\dot{\mathbf x}(t_0)$.
  \thmsource{\cite[ch.~1]{arnold}}
\end{theorem}


\begin{remark}
  Experience shows that the state at a moment also needs \emph{velocities}:
  a stone dropped and a stone thrown from the same point move differently.
  Experience also shows that velocities are \emph{enough}: we do not need to
  know accelerations in advance. That is why the equation has second
  order.
\end{remark}
